\appendix

\section{Proof Details}
\label{sec:proof-details}

\subsection{Finite gradient control on the static ball}
\label{app:step001-gradient}

\begin{lemma}[Finite-ball forward envelope]
\label{lem:p1-i3-step001-forward}
Under Assumptions~\ref{assump:fixed-source-witnesses} and
\ref{assump:robust-tube}, fix a deterministic
$\theta^{(0)}\in\mathbb R^S$ and the finite radius $r>0$.  Define
\[
B:=\max\{1,\|\theta^{(0)}\|_\infty+r\},
\qquad a_0:=1,
\qquad a_\ell:=n_{\ell-1}B a_{\ell-1}
\quad(1\leq\ell<L).
\tag{21}
\]
Then $B,a_0,\ldots,a_{L-1}$ are finite and, for every
$\theta\in B_\infty(\theta^{(0)},r)$, $x\in\mathcal X$, hidden layer
$1\leq\ell<L$, and $j\in[n_\ell]$,
\[
|u_{\ell,j}(\theta,x)|\leq a_\ell,
\qquad
|z_{\ell,j}(\theta,x)|\leq a_\ell.
\tag{22}
\]
For $L=1$ the hidden-layer assertion is empty and $a_0=1$.
\end{lemma}

\begin{proof}
All coordinates of the deterministic vector $\theta^{(0)}$ are finite and
the radius is finite, hence $B<\infty$.  If $\theta$ lies in the closed ball,
then every weight coordinate satisfies
\[
|\theta_{\ell,jk}|
\leq |(\theta^{(0)}_\ell)_{jk}|
      +|\theta_{\ell,jk}-(\theta^{(0)}_\ell)_{jk}|
\leq \|\theta^{(0)}\|_\infty+r\leq B.
\tag{23}
\]
The weak inequalities include the boundary $\|\theta-\theta^{(0)}\|_\infty=r$.
At the input, $z_0(\theta,x)=x$ and every coordinate has magnitude
$1=a_0$.  If all coordinates of $z_{\ell-1}$ have magnitude at most
$a_{\ell-1}$, then
\[
\begin{aligned}
|u_{\ell,j}(\theta,x)|
&=\left|\sum_{k=1}^{n_{\ell-1}}
       \theta_{\ell,jk}z_{\ell-1,k}(\theta,x)\right|\\
&\leq\sum_{k=1}^{n_{\ell-1}}
       |\theta_{\ell,jk}|\,|z_{\ell-1,k}(\theta,x)|\\
&\leq n_{\ell-1}B a_{\ell-1}=a_\ell.
\end{aligned}
\tag{24}
\]
Since $0\leq\sigma(v)\leq|v|$ for every real $v$, the same bound holds
for $z_{\ell,j}$.  Finite induction proves (22), and the recurrence also
gives
\[
a_\ell=B^\ell\prod_{q=0}^{\ell-1}n_q,
\tag{25}
\]
with the empty product equal to one.  Thus all envelopes are finite.
\end{proof}

\begin{proposition}[Exact fixed-selector gradient envelope]
\label{prop:p1-i3-step001-gradient}
Under Assumptions~\ref{assump:fixed-source-witnesses} and
\ref{assump:robust-tube}, and Lemma~\ref{lem:p1-i3-step001-forward}, define
\[
b_L:=1,
\qquad b_\ell:=n_{\ell+1}B b_{\ell+1}\quad(1\leq\ell<L),
\qquad
C_r(\theta^{(0)}):=\max_{1\leq\ell\leq L}b_\ell a_{\ell-1}.
\tag{26}
\]
Then on the full domain defining $G_r(\theta^{(0)})$,
\[
\left|
[\nabla_\theta^{(\kappa)}\ell(yf_\theta(x))]_{\ell,jk}
\right|
\leq b_\ell a_{\ell-1}\leq C_r(\theta^{(0)}),
\tag{27}
\]
for every layer and coordinate.  Consequently
\[
G_r(\theta^{(0)})\leq C_r(\theta^{(0)})<\infty.
\tag{28}
\]
When $L=1$, $C_r(\theta^{(0)})=1$ and $G_r(\theta^{(0)})\leq1$.
\end{proposition}

\begin{proof}
Fix $\theta,x,y$ in the domains of (27).  The forward equations are finite
sums and products, so $f_\theta(x)$ is finite.  Direct differentiation of
$\ell(a)=\log(1+e^{-a})$ gives
\[
\ell'(a)=-\frac1{1+e^a},
\qquad
\left|\frac{\partial}{\partial f}\ell(yf)\right|
=\frac1{1+e^{yf}}\leq1.
\tag{29}
\]
Because $n_L=1$, define the selected output adjoint
\[
\delta_{L,1}:=\frac{\partial}{\partial f}\ell(yf)\bigg|_{f=f_\theta(x)}
=-\frac{y}{1+e^{yf_\theta(x)}}.
\tag{30}
\]
Thus $|\delta_{L,1}|\leq1=b_L$.  The exact protocol-selected recurrence at
the hidden layers is
\[
\delta_{\ell,j}:=\rho_\kappa(u_{\ell,j}(\theta,x))
\sum_{q=1}^{n_{\ell+1}}
\theta_{\ell+1,qj}\delta_{\ell+1,q},
\qquad \ell=L-1,\ldots,1.
\tag{31}
\]
It is valid as the definition of the selected back-propagation value even at
a kink.  Since $0\leq\rho_\kappa(v)\leq1$ for every $v$, including
$\rho_\kappa(0)=\kappa$, the induction hypothesis and (23) give
\[
|\delta_{\ell,j}|
\leq\sum_{q=1}^{n_{\ell+1}}
|\theta_{\ell+1,qj}|\,|\delta_{\ell+1,q}|
\leq n_{\ell+1}B b_{\ell+1}=b_\ell.
\tag{32}
\]
Backward induction yields
\[
b_\ell=B^{L-\ell}\prod_{q=\ell+1}^{L}n_q,
\tag{33}
\]
where $n_L=1$.  The layerwise selected-gradient formula is
\[
[\nabla_\theta^{(\kappa)}\ell(yf_\theta(x))]_{\ell,jk}
=\delta_{\ell,j}z_{\ell-1,k}(\theta,x).
\tag{34}
\]
Combining (32), (34), and Lemma~\ref{lem:p1-i3-step001-forward} proves
(27), and taking the supremum over the domain in (9) proves (28).  The
finiteness of (25), (33), and the finite layer maximum in (26) gives
$C_r<\infty$.  The estimate includes hidden preactivation zeros for every
$\kappa\in[0,1]$ and does not use continuity or attainment of the selected
gradient supremum.

If $L=1$, there are no hidden adjoints and (34) reduces to
\[
\left|[\nabla_\theta^{(\kappa)}\ell(yf_\theta(x))]_{1,1k}\right|
=\frac{|x_k|}{1+e^{yf_\theta(x)}}\leq1,
\tag{35}
\]
so $a_0=b_1=C_r=1$, including zero weights and closed-ball boundary points.
\end{proof}

The two results prove the finite static envelope before any trajectory
membership or gate statement is invoked.  The bound is deterministic for each
initialization and uniform in inputs and labels.

\subsection{All-history containment in the closed tube}
\label{app:step002-containment}

Fix $\theta^{(0)}\in E_r$ and an arbitrary deterministic labeled history
\[
\omega=((x^{(0)},y^{(0)}),\ldots,(x^{(T-1)},y^{(T-1)}))
\in(\mathcal X\times\{-1,+1\})^T.
\]
Let the exact history-driven states satisfy
\[
\theta^{(t+1)}=\theta^{(t)}-\eta g_t,
\qquad
g_t:=\nabla_\theta^{(\kappa)}
\ell(y^{(t)}f_{\theta^{(t)}}(x^{(t)})),
\qquad0\leq t<T,
\tag{36}
\]
and put
\[
D_t:=\|\theta^{(t)}-\theta^{(0)}\|_\infty,
\qquad G:=G_r(\theta^{(0)}).
\tag{37}
\]

\begin{lemma}[One-step closed-ball recurrence]
\label{lem:p1-i3-step002-one-step}
Under Assumptions~\ref{assump:fixed-source-witnesses} and
\ref{assump:robust-tube}, Proposition~\ref{prop:p1-i3-step001-gradient},
and the definitions (36)--(37), if $0\leq t<T$ and $D_t\leq r$, then
\[
\|g_t\|_\infty\leq G,
\qquad D_{t+1}\leq D_t+\eta G.
\tag{38}
\]
The conclusion includes the case $D_t=r$.
\end{lemma}

\begin{proof}
The condition $D_t\leq r$ is exactly
$\theta^{(t)}\in B_\infty(\theta^{(0)},r)$.  Since the history is arbitrary,
$x^{(t)}\in\mathcal X$ and $y^{(t)}\in\{-1,+1\}$.  Therefore the triple
$(\theta^{(t)},x^{(t)},y^{(t)})$ lies in the defining domain of $G_r$, and
\[
\|g_t\|_\infty
=\left\|\nabla_\theta^{(\kappa)}
\ell(y^{(t)}f_{\theta^{(t)}}(x^{(t)}))\right\|_\infty\leq G.
\tag{39}
\]
Using (36), the triangle inequality, $\eta>0$, and (39),
\[
\begin{aligned}
D_{t+1}
&=\|\theta^{(t)}-\theta^{(0)}-\eta g_t\|_\infty\\
&\leq\|\theta^{(t)}-\theta^{(0)}\|_\infty+\eta\|g_t\|_\infty\\
&\leq D_t+\eta G.
\end{aligned}
\tag{40}
\]
No sign, descent, or cancellation property is used.
\end{proof}

\begin{proposition}[All-history closed-tube containment]
\label{prop:p1-i3-step002-closed-tube}
Under Assumptions~\ref{assump:fixed-source-witnesses} and
\ref{assump:robust-tube}, Proposition~\ref{prop:p1-i3-step001-gradient},
and Lemma~\ref{lem:p1-i3-step002-one-step}, fix any initialization
$\theta^{(0)}\in E_r$ and any deterministic labeled history
\[
\omega=((x^{(0)},y^{(0)}),\ldots,(x^{(T-1)},y^{(T-1)}))
\in(\mathcal X\times\{-1,+1\})^T.
\]
Let $\theta^{(t)}$ be generated from $\theta^{(0)}$ by the exact recursion
\[
\theta^{(t+1)}=\theta^{(t)}-\eta\nabla_\theta^{(\kappa)}
\ell\!\left(y^{(t)}f_{\theta^{(t)}}(x^{(t)})\right),
\qquad 0\leq t<T.
\]
Then, for every integer $0\leq t\leq T$,
\[
\|\theta^{(t)}-\theta^{(0)}\|_\infty
\leq t\eta G_r(\theta^{(0)})\leq r.
\tag{41}
\]
\end{proposition}

\begin{proof}
The accepted static envelope gives $0\leq G<\infty$.  Membership in $E_r$
gives only the scalar budget
\[
\eta T G\leq r.
\tag{42}
\]
At the entry state,
\[
D_0=\|\theta^{(0)}-\theta^{(0)}\|_\infty=0=0\cdot\eta G\leq r.
\tag{43}
\]
Suppose $0\leq t<T$ and $D_t\leq t\eta G$.  Since $t\leq T$, (42) gives
\[
D_t\leq t\eta G\leq T\eta G\leq r,
\tag{44}
\]
so the one-step lemma may be applied only after current membership has been
established.  It yields
\[
D_{t+1}\leq D_t+\eta G\leq(t+1)\eta G,
\tag{45}
\]
and $t+1\leq T$ gives
\[
(t+1)\eta G\leq T\eta G\leq r.
\tag{46}
\]
Finite induction proves (41).  The first transition is
$D_1\leq\eta G\leq r$ because $T\geq1$; if $T=1$, (43) and this bound
complete the argument.  If $G=0$, every inequality forces $D_t=0$.  Equality
at the terminal radius is allowed because the ball is closed, and arbitrary
outward increments are already charged by the full sum $t\eta G$.
\end{proof}

Thus the static event supplies no trajectory containment statement by itself:
(41) is the derived all-history containment result that will be used below.  It applies in
particular to source-consistent histories with $y^{(t)}=h(x^{(t)})$.

\subsection{Gate constancy from the positive margin}
\label{app:step003-gates}

\begin{lemma}[Continuous exact preactivation paths]
\label{lem:p1-i3-step003-continuous-preactivation}
Under Assumptions~\ref{assump:fixed-source-witnesses} and
\ref{assump:robust-tube}, fix $\theta^{(0)}$, $r>0$,
$\theta\in B_\infty(\theta^{(0)},r)$, and $x\in\mathcal X$.  For $L\geq2$
define
\[
\gamma(s):=\theta^{(0)}+s(\theta-\theta^{(0)}),
\qquad0\leq s\leq1.
\tag{47}
\]
Then $\gamma(s)$ remains in the closed ball and each exact scalar path
$q_{\ell,j}(s):=u_{\ell,j}(\gamma(s),x)$, $1\leq\ell<L$, is continuous.
More explicitly,
\[
q_{1,j}(s)=\sum_{k=1}^{n_0}\gamma_{1,jk}(s)x_k,
\tag{48}
\]
and, for $2\leq\ell<L$,
\[
q_{\ell,j}(s)=\sum_{k=1}^{n_{\ell-1}}
\gamma_{\ell,jk}(s)\sigma(q_{\ell-1,k}(s)).
\tag{49}
\]
For $L=1$ there are no such paths.
\end{lemma}

\begin{proof}
For every vectorized coordinate,
\[
|\gamma_a(s)-\theta^{(0)}_a|
=s|\theta_a-\theta^{(0)}_a|\leq sr\leq r,
\tag{50}
\]
so the segment, including its endpoints and boundary points, lies in the
closed ball.  Its coordinates are affine and hence continuous.  Formula
(48) is a finite sum of continuous functions.  If all paths at layer
$\ell-1$ are continuous, then coordinatewise ReLU is continuous, and (49) is
a finite sum of products of continuous functions.  Induction over the finite
hidden layers proves continuity.  These are the exact forward equations
evaluated on $\gamma(s)$, not a linearization, and no derivative convention is
used.
\end{proof}

\begin{proposition}[Strict gate constancy on the positive-margin ball]
\label{prop:p1-i3-step003-ball-gates}
Under Assumptions~\ref{assump:fixed-source-witnesses} and
\ref{assump:robust-tube}, Lemma~\ref{lem:p1-i3-step003-continuous-preactivation},
and the local condition $\theta^{(0)}\in E_r$, let
$m=M_r(\theta^{(0)})$.  If $L\geq2$, then $m>0$ and, for every point in the
closed ball, input, hidden layer, and coordinate,
\[
|u_{\ell,j}(\theta,x)|\geq m,
\qquad
\operatorname{sgn}(u_{\ell,j}(\theta,x))
=\operatorname{sgn}(u_{\ell,j}(\theta^{(0)},x)),
\tag{51}
\]
where ordinary $\operatorname{sgn}$ is used only on nonzero reals.  Hence
\[
\mathbf1\{u_{\ell,j}(\theta,x)>0\}
=\mathbf1\{u_{\ell,j}(\theta^{(0)},x)>0\}.
\tag{52}
\]
For $L=1$ the assertion is empty and $M_r=+\infty$ by definition.
\end{proposition}

\begin{proof}
When $L\geq2$, event membership gives $m>0$.  By the infimum definition of
$M_r$, every point $\vartheta$ in the ball satisfies
\[
|u_{\ell,j}(\vartheta,x)|\geq m
\tag{53}
\]
for every hidden coordinate and input.  Along the segment (47), (53) gives
$|q_{\ell,j}(s)|\geq m>0$.  If the endpoint signs differed, continuity from
Lemma~\ref{lem:p1-i3-step003-continuous-preactivation} and the intermediate
value theorem would give an $s_\star\in(0,1)$ with
$q_{\ell,j}(s_\star)=0$, contradicting (53).  Thus the endpoint signs agree;
since both are nonzero, (52) follows.

The argument includes $\theta=\theta^{(0)}$, all ball-boundary points, and
zero parameter coordinates.  A zero preactivation anywhere in the ball would
make $M_r=0$ and hence place the initialization outside $E_r$; that case is
not silently excluded from the later unconditional law.  The proof never
evaluates $\rho_\kappa(0)$ and is independent of $\kappa$.
\end{proof}

\begin{proposition}[All-history initialization gate signature]
\label{prop:p1-i3-step003-history-gates}
Under Assumptions~\ref{assump:fixed-source-witnesses} and
\ref{assump:robust-tube}, Proposition~\ref{prop:p1-i3-step002-closed-tube},
and Proposition~\ref{prop:p1-i3-step003-ball-gates}, fix any initialization
$\theta^{(0)}\in E_r$ and any deterministic labeled history
\[
\omega=((x^{(0)},y^{(0)}),\ldots,(x^{(T-1)},y^{(T-1)}))
\in(\mathcal X\times\{-1,+1\})^T.
\]
Let $\theta^{(t)}$ be generated from $\theta^{(0)}$ by
\[
\theta^{(t+1)}=\theta^{(t)}-\eta\nabla_\theta^{(\kappa)}
\ell\!\left(y^{(t)}f_{\theta^{(t)}}(x^{(t)})\right),
\qquad 0\leq t<T.
\]
If $L\geq2$, then
\[
|u_{\ell,j}(\theta^{(t)},x)|
\geq M_r(\theta^{(0)})>0,
\qquad
\mathbf1\{u_{\ell,j}(\theta^{(t)},x)>0\}
=\mathbf1\{u_{\ell,j}(\theta^{(0)},x)>0\}
\tag{54}
\]
for all $x$, hidden $(\ell,j)$, and $0\leq t\leq T$.  For $L=1$ the
universal statement is empty.
\end{proposition}

\begin{proof}
Proposition~\ref{prop:p1-i3-step002-closed-tube} gives
\[
\|\theta^{(t)}-\theta^{(0)}\|_\infty
\leq t\eta G_r(\theta^{(0)})\leq r
\tag{55}
\]
for every state, including the first and terminal states and equality at the
boundary.  Proposition~\ref{prop:p1-i3-step003-ball-gates} therefore applies
to each $\theta^{(t)}$, every input, and every hidden coordinate, proving
(54).  The history was arbitrary, so the conclusion is all-history rather
than almost-sure.  No gate stability was assumed before (55) was established.
\end{proof}
\subsection{Exact path expansion and aggregate coefficients}
\label{app:step004-paths}

\begin{lemma}[Exact path expansion under the initialization gate signature]
\label{lem:p1-i3-step004-path-expansion}
Under Assumptions~\ref{assump:fixed-source-witnesses} and
\ref{assump:robust-tube}, Proposition~\ref{prop:p1-i3-step003-ball-gates},
and $\theta^{(0)}\in E_r$, fix
$\theta\in B_\infty(\theta^{(0)},r)$ and $x\in\mathcal X$.  For
$p=(i_0,\ldots,i_{L-1})\in\mathcal I_{\rm path}$ define
\[
c_p(\theta):=(\theta_L)_{1,i_{L-1}}
\prod_{\ell=1}^{L-1}(\theta_\ell)_{i_\ell,i_{\ell-1}},
\qquad
\varphi_{\theta^{(0)},p}(x):=x_{i_0}
\prod_{\ell=1}^{L-1}
\mathbf1\{u_{\ell,i_\ell}(\theta^{(0)},x)>0\},
\tag{56}
\]
and let $c(\theta)=(c_p(\theta))_{p\in\mathcal I_{\rm path}}$.  Then
\[
f_\theta(x)=\sum_{p\in\mathcal I_{\rm path}}c_p(\theta)
\varphi_{\theta^{(0)},p}(x)
=\langle c(\theta),\varphi_{\theta^{(0)}}(x)\rangle.
\tag{57}
\]
Empty products are one, so this includes $L=1$.
\end{lemma}

\begin{proof}
Assume first $L\geq2$ and abbreviate
\[
g_{\ell,j}^{(0)}(x):=\mathbf1\{u_{\ell,j}(\theta^{(0)},x)>0\}.
\tag{58}
\]
The ball gate proposition gives
\[
\mathbf1\{u_{\ell,j}(\theta,x)>0\}=g_{\ell,j}^{(0)}(x),
\qquad |u_{\ell,j}(\theta,x)|>0.
\tag{59}
\]
Therefore the exact ReLU equation is
\[
z_{\ell,j}(\theta,x)=g_{\ell,j}^{(0)}(x)
\sum_{k=1}^{n_{\ell-1}}(\theta_\ell)_{j,k}
z_{\ell-1,k}(\theta,x).
\tag{60}
\]
Define finite partial path sums by
\[
P_{0,i_0}(\theta,x):=x_{i_0},
\tag{61}
\]
and, for $1\leq\ell<L$,
\[
P_{\ell,i_\ell}(\theta,x):=g_{\ell,i_\ell}^{(0)}(x)
\sum_{i_{\ell-1}=1}^{n_{\ell-1}}
(\theta_\ell)_{i_\ell,i_{\ell-1}}
P_{\ell-1,i_{\ell-1}}(\theta,x).
\tag{62}
\]
We prove simultaneously $z_{\ell,j}=P_{\ell,j}$ and, for every hidden
layer, the explicit expansion
\[
P_{\ell,i_\ell}(\theta,x)=
\sum_{i_0=1}^{n_0}\cdots\sum_{i_{\ell-1}=1}^{n_{\ell-1}}
x_{i_0}
\left(\prod_{q=1}^{\ell}g_{q,i_q}^{(0)}(x)\right)
\left(\prod_{q=1}^{\ell}(\theta_q)_{i_q,i_{q-1}}\right).
\tag{63}
\]
At layer zero, $z_{0,i_0}=x_{i_0}=P_{0,i_0}$.  If the assertion holds at
layer $\ell-1$, substituting it into (60) gives (62), and distributing the
finite sum appends exactly the new factor
$g_{\ell,i_\ell}^{(0)}(x)(\theta_\ell)_{i_\ell,i_{\ell-1}}$, proving (63).
No division or nonzero weight condition is used.

At the scalar output, $n_L=1$ and
\[
\begin{aligned}
f_\theta(x)
&=\sum_{i_0=1}^{n_0}\cdots\sum_{i_{L-1}=1}^{n_{L-1}}
x_{i_0}\left(\prod_{q=1}^{L-1}g_{q,i_q}^{(0)}(x)\right)
\left[(\theta_L)_{1,i_{L-1}}
\prod_{q=1}^{L-1}(\theta_q)_{i_q,i_{q-1}}\right]\\
&=\sum_{p\in\mathcal I_{\rm path}}c_p(\theta)
\varphi_{\theta^{(0)},p}(x).
\end{aligned}
\tag{64}
\]
If $L=1$, then $\mathcal I_{\rm path}=[n_0]$ and the original output equation
is already
\[
f_\theta(x)=\sum_{i_0=1}^{n_0}(\theta_1)_{1,i_0}x_{i_0}
=\sum_{i_0\in[n_0]}c_{i_0}(\theta)
\varphi_{\theta^{(0)},i_0}(x).
\tag{65}
\]
Zero or canceling products simply give zero summands or exact cancellation;
the finite equality remains valid.
\end{proof}

\begin{proposition}[Exact latter-half aggregate coefficient and tie identity]
\label{prop:p1-i3-step004-aggregate}
Under Assumptions~\ref{assump:fixed-source-witnesses} and
\ref{assump:robust-tube}, Propositions~\ref{prop:p1-i3-step002-closed-tube}
and~\ref{prop:p1-i3-step003-history-gates}, and
Lemma~\ref{lem:p1-i3-step004-path-expansion}, fix an initialization in $E_r$
and a deterministic labeled history with its exact states.  Set
\[
J_T:=\{\lceil T/2\rceil,\ldots,T\},
\qquad
w_\omega:=\sum_{t\in J_T}c(\theta^{(t)}(\omega)).
\tag{66}
\]
Then, for every $x\in\mathcal X$,
\[
A_\omega(x):=\sum_{t\in J_T}f_{\theta^{(t)}(\omega)}(x)
=\langle w_\omega,\varphi_{\theta^{(0)}}(x)\rangle.
\tag{67}
\]
For a source-consistent history this is the setting score, and applying
$\operatorname{sign}_{s_0}$ gives the same prediction, including at a zero
score.
\end{proposition}

\begin{proof}
Proposition~\ref{prop:p1-i3-step002-closed-tube} places each state in the
closed ball, and Proposition~\ref{prop:p1-i3-step003-history-gates} supplies
the initialization gate signature at every state.  Lemma
\ref{lem:p1-i3-step004-path-expansion} therefore gives, for each $t\in J_T$,
\[
f_{\theta^{(t)}(\omega)}(x)
=\left\langle c(\theta^{(t)}(\omega)),
\varphi_{\theta^{(0)}}(x)\right\rangle.
\tag{68}
\]
The set $J_T$ is finite and nonempty because $T\geq1$.  Summing (68) and
using bilinearity proves (67).  On a source-consistent history the left side
is exactly (5), so equal real scores have equal
$\operatorname{sign}_{s_0}$ values and identical strict-error indicators.
When $T=1$, $J_T=\{1\}$.  Zero coefficients, cancellations, and a common score
of zero do not change the equality; at zero both predictors use $s_0$.
\end{proof}

\subsection{The conditional infimum comparison}
\label{app:step005-risk}

Fix $\mathcal D\in\Delta(\mathcal X)$, $h\in\mathcal H$, and a stable
initialization $\theta^{(0)}\in E_r$.  For a training tuple
$\mathbf x=(x^{(0)},\ldots,x^{(T-1)})\in\mathcal X^T$, use the source-consistent
history
\[
\omega_h(\mathbf x):=((x^{(0)},h(x^{(0)})),\ldots,
(x^{(T-1)},h(x^{(T-1)}))).
\tag{69}
\]
Let $w_{\mathbf x}$ be the coefficient in Proposition
\ref{prop:p1-i3-step004-aggregate}, and abbreviate the corresponding score and
predictor by
\[
A_{\theta^{(0)},\mathbf x}(x):=
\sum_{t=\lceil T/2\rceil}^{T}
f_{\theta^{(t)}(\theta^{(0)},\omega_h(\mathbf x))}(x),
\qquad
\widehat h_{\theta^{(0)},\mathbf x}(x):=
\operatorname{sign}_{s_0}(A_{\theta^{(0)},\mathbf x}(x)).
\tag{70}
\]

\begin{lemma}[Pathwise exact risk identity]
\label{lem:p1-i3-step005-pathwise-risk}
Under Assumptions~\ref{assump:fixed-source-witnesses} and
\ref{assump:robust-tube}, Proposition~\ref{prop:p1-i3-step004-aggregate},
and fixed $\mathcal D,h,\theta^{(0)}\in E_r$, every
$\mathbf x\in\mathcal X^T$ has a vector $w_{\mathbf x}\in\mathbb R^{d_{\rm path}}$
such that
\[
\operatorname{sign}_{s_0}(A_{\theta^{(0)},\mathbf x}(x))
=\operatorname{sign}_{s_0}
\left(\langle w_{\mathbf x},\varphi_{\theta^{(0)}}(x)\rangle\right)
\tag{71}
\]
for every $x$, and
\[
R_{\mathcal D,h}(w_{\mathbf x},\varphi_{\theta^{(0)}})
=\mathcal L_{\mathcal D,h}(\widehat h_{\theta^{(0)},\mathbf x}).
\tag{72}
\]
\end{lemma}

\begin{proof}
The history (69) is deterministic and source-consistent.  Proposition
\ref{prop:p1-i3-step004-aggregate} gives the real-score equality
\[
A_{\theta^{(0)},\mathbf x}(x)
=\langle w_{\mathbf x},\varphi_{\theta^{(0)}}(x)\rangle
\tag{73}
\]
for every evaluation input.  Applying the same tie-resolved sign, multiplying
by the same $h(x)$, and testing the same strict inequality gives pointwise
equality of the error indicators, hence (72).  The conclusion remains exact
when the common score is zero, and no coefficient nonzero or margin condition
is needed.  For $T=1$, the index set contains the single terminal state.
\end{proof}

\begin{proposition}[Conditional infimum-before-expectation comparison]
\label{prop:p1-i3-step005-conditional-comparison}
Under Assumptions~\ref{assump:fixed-source-witnesses} and
\ref{assump:robust-tube}, Proposition~\ref{prop:p1-i3-step004-aggregate},
and Lemma~\ref{lem:p1-i3-step005-pathwise-risk}, fix
$\mathcal D\in\Delta(\mathcal X)$, $h\in\mathcal H$, and
$\theta^{(0)}\in E_r$ before drawing $\mathbf X\sim\mathcal D^T$.  Define
\[
F_{\mathcal D,h}(\theta^{(0)}):=
\inf_{w\in\mathbb R^{d_{\rm path}}}
R_{\mathcal D,h}(w,\varphi_{\theta^{(0)}}).
\tag{74}
\]
Then
\[
0\leq F_{\mathcal D,h}(\theta^{(0)})\leq1,
\qquad
F_{\mathcal D,h}(\theta^{(0)})
\leq\mathbb E_{\mathbf X\sim\mathcal D^T}
\left[\mathcal L_{\mathcal D,h}
(\widehat h_{\theta^{(0)},\mathbf X})\right].
\tag{75}
\]
The infimum need not be attained.
\end{proposition}

\begin{proof}
Every risk in (74) is a probability in $[0,1]$, and the coefficient space is
nonempty, so its infimum is a finite value in $[0,1]$.  For every fixed
$\mathbf x$, the pathwise lemma supplies the feasible vector $w_{\mathbf x}$;
the defining lower-bound property of an infimum gives
\[
F_{\mathcal D,h}(\theta^{(0)})
\leq R_{\mathcal D,h}(w_{\mathbf x},\varphi_{\theta^{(0)}})
=\mathcal L_{\mathcal D,h}
(\widehat h_{\theta^{(0)},\mathbf x}).
\tag{76}
\]
The left side is fixed before the tuple is drawn.  Multiplying (76) by the
nonnegative mass $\mathcal D^T(\mathbf x)$ and summing over the finite set
$\mathcal X^T$ gives
\[
\begin{aligned}
\mathbb E_{\mathbf X\sim\mathcal D^T}
\left[\mathcal L_{\mathcal D,h}
(\widehat h_{\theta^{(0)},\mathbf X})\right]
&=\sum_{\mathbf x\in\mathcal X^T}\mathcal D^T(\mathbf x)
\mathcal L_{\mathcal D,h}(\widehat h_{\theta^{(0)},\mathbf x})\\
&\geq\sum_{\mathbf x\in\mathcal X^T}\mathcal D^T(\mathbf x)
F_{\mathcal D,h}(\theta^{(0)})\\
&=F_{\mathcal D,h}(\theta^{(0)}).
\end{aligned}
\tag{77}
\]
This is pointwise feasible-candidate comparison followed by averaging, not an
interchange of an infimum with an expectation, and it does not require a
measurable minimizing coefficient.
\end{proof}


\subsection{Unconditional gate-law witness}
\label{app:step006-witness}

Let $\Theta=\mathbb R^S$ and let $\gamma$ be the product Gaussian
initialization law in (3).  Define
\[
\Phi:\Theta\longrightarrow\{\text{maps }\mathcal X\to\mathbb R^{d_{\rm path}}\},
\qquad \Phi(\vartheta):=\varphi_\vartheta,
\tag{78}
\]
and, for fixed $\mathcal D,h$,
\[
F_{\mathcal D,h}(\vartheta):=
\inf_{w\in\mathbb R^{d_{\rm path}}}
R_{\mathcal D,h}(w,\Phi(\vartheta)).
\tag{79}
\]
This definition applies on both the stable event and its complement.

\begin{lemma}[Finite gate law and pushforward identity]
\label{lem:p1-i3-step006-gate-pushforward}
Under Assumption~\ref{assump:fixed-source-witnesses}, the map $\Phi$ is Borel
measurable and has finite range $\Gamma=\Phi(\Theta)$, with
\[
|\Gamma|\leq3^{d_{\rm path}|\mathcal X|}.
\tag{80}
\]
For every fixed $\mathcal D\in\Delta(\mathcal X)$ and $h\in\mathcal H$,
$F_{\mathcal D,h}$ is Borel, takes values in $[0,1]$, and
\[
\mathbb E_{\varphi\sim\mathcal P_{\rm gate}}
\left[\inf_{w\in\mathbb R^{d_{\rm path}}}
R_{\mathcal D,h}(w,\varphi)\right]
=\mathbb E_{\theta^{(0)}\sim\gamma}
[F_{\mathcal D,h}(\theta^{(0)})].
\tag{81}
\]
No optimizing coefficient is selected measurably.
\end{lemma}

\begin{proof}
For fixed $x$, every initialization preactivation is continuous in
$\vartheta$ by finite matrix multiplication and continuous coordinatewise
ReLU.  Hence each strict indicator
$\mathbf1\{u_{\ell,j}(\vartheta,x)>0\}$ is Borel.  Each coordinate of $\Phi$
is a finite product of such indicators and one input coordinate, and therefore
belongs to $\{-1,0,+1\}$.  There are only
$d_{\rm path}|\mathcal X|$ coordinates, proving measurability and (80).
For $L=1$, the empty product gives $\Phi(\vartheta)(x)=x$ and the range is a
singleton.

For $\psi\in\Gamma$, define the finite-table function
\[
\Psi_{\mathcal D,h}(\psi):=
\inf_{w\in\mathbb R^{d_{\rm path}}}R_{\mathcal D,h}(w,\psi).
\tag{82}
\]
Every risk is in $[0,1]$ and the coefficient space is nonempty, so the table
is finite-valued even if its infimum is not attained.  Equation (79) is
$F_{\mathcal D,h}=\Psi_{\mathcal D,h}\circ\Phi$, hence $F$ is Borel.  Since
$\mathcal P_{\rm gate}=\Phi_\#\gamma$, summing over the finite range gives
\[
\begin{aligned}
\mathbb E_{\varphi\sim\mathcal P_{\rm gate}}
\left[\inf_wR_{\mathcal D,h}(w,\varphi)\right]
&=\sum_{\psi\in\Gamma}\Psi_{\mathcal D,h}(\psi)
\gamma(\Phi^{-1}(\{\psi\}))\\
&=\int_\Theta\Psi_{\mathcal D,h}(\Phi(\vartheta))\,d\gamma(\vartheta)\\
&=\mathbb E_{\theta^{(0)}\sim\gamma}[F_{\mathcal D,h}(\theta^{(0)})].
\end{aligned}
\tag{83}
\]
Initializations inducing the same map are combined by the pushforward mass;
no coefficient selector appears.
\end{proof}

\begin{proposition}[Unconditional event-split risk bound]
\label{prop:p1-i3-step006-event-split}
Under Assumptions~\ref{assump:fixed-source-witnesses},
\ref{assump:universal-expected-success}, and~\ref{assump:robust-tube}, and
Proposition~\ref{prop:p1-i3-step005-conditional-comparison}, for every fixed
$\mathcal D\in\Delta(\mathcal X)$ and $h\in\mathcal H$,
\[
\mathbb E_{\theta^{(0)}\sim\gamma}
[F_{\mathcal D,h}(\theta^{(0)})]\leq\varepsilon+\delta_0.
\tag{84}
\]
The expectation is under the unconditional Gaussian law.
\end{proposition}

\begin{proof}
For a fixed pair $(x,y)\in\mathcal X\times\{-1,+1\}$, define the full
one-step update map
\[
\mathsf T_{x,y}(\theta)
:=\theta-\eta\nabla_\theta^{(\kappa)}
\ell\!\left(yf_\theta(x)\right),
\qquad \theta\in\Theta.
\]
This map is Borel.  Indeed, every forward preactivation and activation is
continuous in $\theta$.  The selected gate factor
$\rho_\kappa(u_{\ell,j}(\theta,x))$ is Borel, including at zero; the output
adjoint $-y/(1+e^{yf_\theta(x)})$ is continuous; and the finite backward
recursion and the coordinate formula
$[\nabla_\theta^{(\kappa)}\ell(yf_\theta(x))]_{\ell,jk}
=\delta_{\ell,j}z_{\ell-1,k}(\theta,x)$ use only finite sums and products.
Thus the selected gradient is a Borel map from $\Theta$ to $\Theta$, and
subtracting it from the identity proves the claim for $\mathsf T_{x,y}$.

Fix a tuple
$\mathbf x=(x^{(0)},\ldots,x^{(T-1)})\in\mathcal X^T$ and define the
generated-state maps recursively by
\[
\mathsf U_{\mathbf x,0}(\vartheta):=\vartheta,
\qquad
\mathsf U_{\mathbf x,t+1}(\vartheta)
:=\mathsf T_{x^{(t)},h(x^{(t)})}
   (\mathsf U_{\mathbf x,t}(\vartheta)),
\qquad 0\leq t<T.
\]
The base map is the identity.  If $\mathsf U_{\mathbf x,t}$ is Borel, then
the displayed composition with the Borel map
$\mathsf T_{x^{(t)},h(x^{(t)})}$ is Borel.  Finite induction therefore proves
that every generated state map $\mathsf U_{\mathbf x,t}:\Theta\to\Theta$ is
Borel for $0\leq t\leq T$.

For each evaluation input $v\in\mathcal X$, the finite-horizon score and
predictor are consequently the Borel compositions
\[
A_{\vartheta,\mathbf x}(v)
:=\sum_{t=\lceil T/2\rceil}^{T}
f_{\mathsf U_{\mathbf x,t}(\vartheta)}(v),
\qquad
\widehat h_{\vartheta,\mathbf x}(v)
:=\operatorname{sign}_{s_0}(A_{\vartheta,\mathbf x}(v)).
\]
The tie map is Borel.  Hence, for an initialization $\vartheta$ and the fixed
tuple $\mathbf x$, the prescribed learner loss is the finite sum
\[
Z_{\mathcal D,h}(\vartheta,\mathbf x):=
\mathcal L_{\mathcal D,h}(\widehat h_{\vartheta,\mathbf x})
=\sum_{v\in\mathcal X}\mathcal D(v)
\mathbf1\!\left\{
\operatorname{sign}_{s_0}(A_{\vartheta,\mathbf x}(v))h(v)<0
\right\}\in[0,1]
\tag{85}
\]
on the source-consistent labeled tuple
$((x^{(t)},h(x^{(t)})))_{t=0}^{T-1}$.  It is Borel in $\vartheta$ for each
fixed tuple.  Since $\mathcal X^T$ is finite with its discrete sigma-algebra,
$Z_{\mathcal D,h}$ is jointly measurable on $\Theta\times\mathcal X^T$.
Define the conditional sample average
\[
g_{\mathcal D,h}(\vartheta):=
\mathbb E_{\mathbf X\sim\mathcal D^T}
[Z_{\mathcal D,h}(\vartheta,\mathbf X)]
=\sum_{\mathbf x\in\mathcal X^T}\mathcal D^T(\mathbf x)
Z_{\mathcal D,h}(\vartheta,\mathbf x),
\tag{86}
\]
which is measurable and lies in $[0,1]$.

For every $\vartheta\in E_r$, the conditional comparison gives
\[
F_{\mathcal D,h}(\vartheta)\leq g_{\mathcal D,h}(\vartheta).
\tag{87}
\]
For every $\vartheta\in E_r^c$, the probability definition gives
\[
0\leq F_{\mathcal D,h}(\vartheta)\leq1.
\tag{88}
\]
Splitting the same unconditional expectation therefore yields
\[
\begin{aligned}
\mathbb E_{\theta^{(0)}}F_{\mathcal D,h}(\theta^{(0)})
&=\mathbb E[\mathbf1_{E_r}F_{\mathcal D,h}]
 +\mathbb E[\mathbf1_{E_r^c}F_{\mathcal D,h}]\\
&\leq\mathbb E[\mathbf1_{E_r}g_{\mathcal D,h}]
 +\Pr(E_r^c).
\end{aligned}
\tag{89}
\]
By independence and the finite sum (86), bounded tower averaging gives
\[
\begin{aligned}
\mathbb E[\mathbf1_{E_r}g_{\mathcal D,h}]
&=\mathbb E_{\theta^{(0)},\mathbf X}
[\mathbf1_{E_r}Z_{\mathcal D,h}(\theta^{(0)},\mathbf X)]\\
&\leq\mathbb E_{\theta^{(0)},\mathbf X}
[Z_{\mathcal D,h}(\theta^{(0)},\mathbf X)]\\
&\leq\varepsilon,
\end{aligned}
\tag{90}
\]
where the last inequality is Assumption~\ref{assump:universal-expected-success}.
Assumption~\ref{assump:robust-tube} gives
\[
\Pr(E_r^c)\leq\delta_0.
\tag{91}
\]
Substitution in (89) proves (84).  The complement is charged once by its
unit risk bound; it is neither discarded nor used to condition the feature
law.
\end{proof}

\begin{proposition}[Universal probabilistic-dimension witness]
\label{prop:p1-i3-step006-dc-witness}
Under Assumptions~\ref{assump:fixed-source-witnesses},
\ref{assump:universal-expected-success}, and~\ref{assump:robust-tube},
together with Lemma~\ref{lem:p1-i3-step006-gate-pushforward} and
Proposition~\ref{prop:p1-i3-step006-event-split}, the one law
$\mathcal P_{\rm gate}$ satisfies, for every
$\mathcal D\in\Delta(\mathcal X)$ and $h\in\mathcal H$,
\[
\mathbb E_{\varphi\sim\mathcal P_{\rm gate}}
\left[\inf_{w\in\mathbb R^{d_{\rm path}}}
R_{\mathcal D,h}(w,\varphi)\right]\leq\varepsilon+\delta_0,
\tag{92}
\]
and hence
\[
\operatorname{dc}_{\varepsilon+\delta_0}(\mathcal H)
\leq d_{\rm path}.
\tag{93}
\]
\end{proposition}

\begin{proof}
The architecture, Gaussian law, and rule $\Phi(\vartheta)=\varphi_\vartheta$
are fixed by Assumption~\ref{assump:fixed-source-witnesses} before any
distribution or target.  Thus their pushforward $\mathcal P_{\rm gate}$ is
one law independent of later choices.  For arbitrary $\mathcal D,h$, (81)
and (84) give (92).  Since the law was fixed before this arbitrary choice, the
quantifier order is $\exists\mathcal P_{\rm gate}\,\forall\mathcal D\,\forall h$.
The definition (1) then gives (93).
\end{proof}

\subsection{Polynomial dimension and the two-epsilon specialization}
\label{app:step007-specialization}

\begin{lemma}[Polynomial path-count bound]
\label{lem:p1-i3-step007-path-count}
Under Assumption~\ref{assump:constant-depth}, positive integer widths
$n_0,\ldots,n_L$ with $n_L=1$, and the setting definitions
$S=\sum_{j=1}^{L}n_jn_{j-1}$ and
$d_{\rm path}=\prod_{j=0}^{L-1}n_j$,
\[
d_{\rm path}\leq S^L\leq S^{L_0}.
\tag{94}
\]
\end{lemma}

\begin{proof}
For each $j\in\{0,\ldots,L-1\}$, positivity of $n_{j+1}$ gives
\[
n_j\leq n_{j+1}n_j\leq\sum_{k=1}^{L}n_kn_{k-1}=S.
\tag{95}
\]
Multiplying the $L$ inequalities yields
\[
d_{\rm path}=\prod_{j=0}^{L-1}n_j\leq S^L.
\tag{96}
\]
Every summand in $S$ is a positive integer, so $S\geq1$.  Since
$L\leq L_0$,
\[
S^{L_0}=S^LS^{L_0-L}\geq S^L.
\tag{97}
\]
If $L=1$, then $n_1=1$ and $S=n_0=d_{\rm path}=S^L$.  If $S=1$, positivity
forces $L=1$ and all three quantities in (94) equal one.
\end{proof}

\begin{proposition}[Public error-threshold specialization]
\label{prop:p1-i3-step007-public-specialization}
Under Assumptions~\ref{assump:fixed-source-witnesses},
\ref{assump:universal-expected-success}, and~\ref{assump:robust-tube}, and
Proposition~\ref{prop:p1-i3-step006-dc-witness}, the same unconditional law
$\mathcal P_{\rm gate}$ satisfies
\[
\forall\mathcal D\in\Delta(\mathcal X)\quad
\forall h\in\mathcal H,
\qquad
\mathbb E_{\varphi\sim\mathcal P_{\rm gate}}
\left[\inf_{w\in\mathbb R^{d_{\rm path}}}
R_{\mathcal D,h}(w,\varphi)\right]\leq2\varepsilon
\tag{98}
\]
and consequently
\[
\operatorname{dc}_{2\varepsilon}(\mathcal H)
\leq\operatorname{dc}_{\varepsilon+\delta_0}(\mathcal H)
\leq d_{\rm path}.
\tag{99}
\]
\end{proposition}

\begin{proof}
The robust-tube range gives $0\leq\delta_0\leq\varepsilon$, hence
\[
\varepsilon+\delta_0\leq2\varepsilon.
\tag{100}
\]
The witness bound (92) therefore proves (98) under the unchanged law and
dimension.  If $0\leq a\leq b$, every law admissible at level $a$ is
admissible at level $b$, so
\[
\operatorname{dc}_b(\mathcal H)\leq\operatorname{dc}_a(\mathcal H).
\tag{101}
\]
Apply this with $a=\varepsilon+\delta_0$ and $b=2\varepsilon$; nonemptiness
of the stricter admissible set is supplied by Proposition
\ref{prop:p1-i3-step006-dc-witness}.  If $\varepsilon=0$, then $\delta_0=0$
and both thresholds are zero; if $\delta_0=\varepsilon$, they coincide.
\end{proof}

\subsection{Proof of the Main Theorem}
\label{app:main-proof}

\begin{proof}[Proof of Theorem~\ref{thm:main}]
Proposition~\ref{prop:p1-i3-step006-dc-witness} proves (18) with one
unconditional law and gives
\[
\operatorname{dc}_{\varepsilon+\delta_0}(\mathcal H)
\leq d_{\rm path}.
\tag{102}
\]
Lemma~\ref{lem:p1-i3-step007-path-count} gives
$d_{\rm path}\leq S^L\leq S^{L_0}$, proving (19).  Finally,
Proposition~\ref{prop:p1-i3-step007-public-specialization} reuses that same
law and gives the first inequality in (20); combining it with (102) and the
path-count bound gives the stated $S^{L_0}$ corollary.  All steps preserve the
fixed tie rule, finite-horizon expectation mode, and initialization-first
quantifier order.
\end{proof}
